\documentclass[a4paper,10pt]{article}
\usepackage[spanish]{babel}
\usepackage{amsmath}
\usepackage{siunitx}
\usepackage{graphicx}
\usepackage{multicol}
\usepackage{geometry}
\usepackage{fancyhdr}
\usepackage{tikz}
\usepackage{eso-pic}
\usepackage{helvet}
\renewcommand{\familydefault}{\sfdefault}


\geometry{top=.5cm, bottom=.5cm, left=1.5cm, right=.5cm}

\pagestyle{fancy}
\fancyhf{}
\cfoot{\thepage}


\begin{document}

\begin{center}
    \Large \textbf{Evaluación A – Electrotecnia II (Capacitores)} \\
    \normalsize Duración: 60 minutos \\
    Alumno/a: \underline{\hspace{6cm}} \hspace{0.5cm} Fecha: \underline{\hspace{2cm}}
\end{center}

\vspace{0.3cm}


\vspace{0.2cm}

\begin{multicols}{2}

\section*{Problemas}

\begin{enumerate}
\item \textbf{Permitividad absoluta.} Calcular \(\varepsilon\) para:
\begin{itemize}
\item nafta \(\left(\varepsilon_r=2{,}35\right)\),
\item agua \(\left(\varepsilon_r=80{,}5\right)\).
\end{itemize}

\item \textbf{Capacitor en vacío.} Dos placas cuadradas de \(10\,\si{cm}\times 10\,\si{cm}\) separadas \(d=0{,}10\,\si{mm}\) en \textit{vacío}. Hallar \(C\).

\item \textbf{Dieléctrico: mica.} Un capacitor plano con \(S=1{,}00\,\si{m^2}\) y \(d=0{,}010\,\si{mm}\) relleno de \textit{mica} \(\left(\varepsilon_r=5{,}6\right)\). Calcular \(C\).

\item \textbf{Comparación dieléctrico.} Para \(S=0{,}10\,\si{m^2}\) y \(d=0{,}0010\,\si{mm}\):
\begin{enumerate}
\item Calcular \(C\) en \textit{vacío}.
\item Recalcular \(C\) si el dieléctrico es \textit{vidrio} \(\left(\varepsilon_r=4{,}7\right)\).
\end{enumerate}

\item \textbf{Carga y energía.} Un capacitor de \(C=10{,}0\,\si{\micro F}\) se conecta a \(V=120\,\si{V}\).
\begin{enumerate}
\item Hallar \(Q\).
\item Hallar la energía almacenada \(U\).
\end{enumerate}

\item \textbf{Área desconocida.} Un capacitor plano tiene \(C=100\,\si{pF}\) y separación \(d=0{,}10\,\si{mm}\). Si \(\varepsilon=5{,}00\times 10^{-11}\,\si{F/m}\), determinar el área \(S\).

\item \textbf{Identificación de material.} Se miden \(S=0{,}0100\,\si{m^2}\), \(d=0{,}200\,\si{mm}\) y \(C=4{,}425\,\si{pF}\).
\begin{enumerate}
\item Hallar \(\varepsilon\).
\item Hallar \(\varepsilon_r\).
\end{enumerate}

\item \textbf{Aplicación directa.} Dos placas de \(500\,\si{cm^2}\) separadas \(0{,}050\,\si{mm}\) con \(\varepsilon=7{,}00\times 10^{-12}\,\si{F/m}\). Calcular \(C\).

\item \textbf{Cambio geométrico y dieléctrico.} Un capacitor plano en \textit{vacío} tiene capacitancia \(C_0\).
Luego se \emph{duplica} la separación entre placas y se inserta un dieléctrico con \(\varepsilon_r=3{,}0\).
\begin{enumerate}
\item Expresar la nueva capacitancia \(C\) en función de \(C_0\).
\item Si \(C_0=120\,\si{pF}\), calcular \(C\).
\end{enumerate}

\end{enumerate}

\end{multicols}

\end{document}
